% Part: first-order-logic
% Chapter: models-theories
% Section: introduction

\documentclass[../../../include/open-logic-section]{subfiles}

\begin{document}

\olfileid{fol}{mat}{int}
\olsection{پېژندنه}

\begin{explain}
د بديهي طريقې وده د ساينس په تاريخ کښې يوه ستره لاسته راوړنه ده،
او د رياضۍ په تاريخ کښې ځانګړے اهميت لري۔ د يوه ډګر بديهي وده ډېرې
پوښتنې روښانوي: دا ډګر د څه په باب دے؟ تر ټولو بنسټيز مفاهيم ئې کوم
دي؟ هغوئ له يو بل سره څنګه تړلي دي؟ ايا د ډګر ټول مفاهيم د دغو
بنسټيزو مفاهيمو له مخې تعريفېدے شي؟ دا مفاهيم د کومو قوانينو پيروي
کوي، او بايد ئې وکړي؟

بديهي طريقه او منطق د يو بل دپاره جوړ شوي دي۔ صوري منطق د بديهي
تيوريو د جوړولو، د تيورۍ له بديهي اصولو څخه په کره ټاکلې طريقه د
قضيو د ثابتولو، او د بديهي اصولو پوره کوونکو ټولو نظامونو د خاصيتونو
د منظمې څېړنې وسيلې برابروي۔
\end{explain}

\begin{defn}
د !!{sentence}s يو سټ~$\Gamma$ هله او يوازې هله \emph{بند} دے چې
هر کله $\Gamma \Entails !A$ وي، نو $!A \in \Gamma$ هم وي۔ د
!!{sentence}s د يوه سټ~$\Gamma$ \emph{بندښت} دا دے:
$\Setabs{!A}{\Gamma \Entails !A}$۔

وايو چې~$\Gamma$ د جملو د يوه سټ~$\Delta$ له خوا \emph{په بديهي
اصولو ټاکل شوے} دے، که~$\Gamma$ د~$\Delta$ بندښت وي۔
\end{defn}

\begin{explain}
يوه بديهي تيوري د هغو جملو سټ ګڼلے شو چې د هغې د بديهي اصولو
سټ~$\Delta$ په بديهي اصولو ټاکلے وي۔ په بل عبارت، کله چې د لومړۍ درجې ژبه د هغه
ساينس د بنسټيزو مفاهيمو دپاره غيرمنطقي سمبولونه ولري چې غواړو په
بديهي ډول ئې وڅېړو، او ورسره د !!{sentence}s داسې سټ ولرو چې د
ساينس بنسټيز قوانين بيانوي، نو تيوري په دې ژبه کښې د ټولو هغو
!!{sentence}s په توګه انځورولے شو چې له بديهي اصولو څخه لازمېږي۔
دا لړۍ له هغو ساده بېلګو څخه، چې يوازې يو بنسټيز مفهوم او ساده
بديهي اصول لري، لکه د جزوي ترتيبونو تيوري، تر نيوټني ميخانيک غوندې
پېچلو تيوريو پورې غځېږي۔

لاندې مهم منطقي حقيقتونه دي چې دغه صوري چلند د بديهي طريقې دپاره
دومره ارزښتمن کوي۔ فرض کړئ~$\Gamma$ د يوې تيورۍ بديهي نظام، يعنې د
جملو يو سټ دے۔
\begin{enumerate}
\item په کره توګه ويلے شو چې يو بديهي نظام د !!{structure}s مطلوبه
  ټولګه کله رانغاړي۔ يعنې که د !!{structure}s يوه ټاکلې ټولګه مو په
  پام کښې وي، بديهي نظام~$\Gamma$ هغه هله او يوازې هله په برياليتوب
  رانغاړي چې مطلوب !!{structure}s کټ مټ هغه~$\Struct M$ وي چې
  $\Sat{M}{\Gamma}$ وي۔
\item کېدے شي په دې کار کښې ځکه پاتې راشو چې داسې~$\Struct M$ وي
  چې $\Sat{M}{\Gamma}$، خو~$\Struct M$ زموږ له مطلوبو
  !!{structure}s څخه نۀ وي۔ دا ښايي دې ته مو اړ کړي چې داسې بديهي
  اصول ورزيات کړو چې په~$\Struct M$ کښې رښتوني نۀ وي۔
\item که لږ تر لږه په دې کښې بريالي يو چې~$\Gamma$ په ټولو مطلوبو
  !!{structure}s کښې رښتونے وي، نو هر کله چې
  $\Gamma \Entails !A$ وي، جمله~$!A$ هم په ټولو مطلوبو
  !!{structure}s کښې رښتونې ده۔ نو د منطقي وسيلو (لکه د
  !!{derivation} طريقو) په کارولو ښودلے شو چې جملې په ټولو مطلوبو
  !!{structure}s کښې رښتونې دي: بس دا وښيو چې له بديهي اصولو څخه
  لازمېږي۔
\item کله کله مو کوم مطلوب !!{structure}s په پام کښې نۀ وي، بلکې
  له خپلو بديهي اصولو پيل کوو: ځينې بنسټيز مفاهيم اخلو چې غواړو
  ټاکلي قوانين پوره کړي، او دغه قوانين په يوه بديهي نظام کښې ثبتوو۔
  سمدستي دا تصديق کول غواړو چې بديهي اصول يو بل نۀ نقضوي؛ که ئې
  وکړي، هېڅ مفاهيم دغه قوانين نۀ شي پوره کولے او تيوري مو ناسازګاره
  جوړه کړې ده۔ د~$\Gamma$ د يوه مدل په موندلو تصديق کولے شو چې داسې
  نۀ دي شوي۔ او که تيوري مو مدلونه ولري، په منطقي طريقو ئې څېړلے او
  مدلونه هم جوړولے شو۔
\item د بديهي اصولو استقلال هم يوه مهمه پوښتنه ده۔ کېدے شي يو بديهي
  اصل په رښتيا د نورو پايله وي او له همدې امله زائد وي۔ ثابتولے شو
  چې په~$\Gamma$ کښې بديهي اصل~$!A$ زائد دے، که
  $\Gamma \setminus \{!A\} \Entails !A$ ثابت کړو۔ دا هم ثابتولے شو
  چې بديهي اصل زائد نۀ دے، که وښيو چې
  $(\Gamma \setminus \{!A\}) \cup \{\lnot !A\}$ د صدق وړ دے۔ د
  بېلګې په توګه، د هندسې د نورو بديهي اصولو څخه د موازي کرښو د اصل
  استقلال همداسې ښودل شوے و۔
\item بله مهمه پوښتنه په يوه تيورۍ کښې د مفاهيمو د تعريفېدنې په باب ده: د
  ژبې ټاکنه دا ټاکي چې د تيورۍ مدلونه له څه جوړ وي۔ خو د تيورۍ هر اړخ
  لازم نۀ دے چې په مدلونو کښې جلا استازيتوب ولري۔ مثلاً هر ترتيب
  $\le$ ورسره سم سخت ترتيب~$<$ ټاکي---چې يو راکړل شي، بل ترې تعريفولے
  شو۔ نو د داسې ترتيب لرونکې تيورۍ په مدل کښې د اړوند سخت ترتيب
  \emph{جلا} شتون لازم نۀ دے۔ په عمومي ډول، يوه اړيکه د نورو له مخې
  کله تعريفېدے شي؟ يوه اړيکه د نورو له مخې کله نۀ شي تعريفېدے (او له
  همدې امله بايد د ژبې په بنسټيزو سمبولونو کښې ورزياته شي)؟
\end{enumerate}
\end{explain}


\end{document}
